% =====================================================================
% VERSION ENSAE, REDUCTION SFE (phase 3, 1er octobre 2026) : remerciements.
% Copie a texte constant de remerciements.tex.
% Ne changent que les titres et les renvois repares, listes dans
% reduction/phase3.md.
% Lu par main_ensae.tex seul : ne jamais le faire lire par main_institut.tex,
% qui garde chapitres/.
% =====================================================================

% =====================================================================
% EXERGUE ET REMERCIEMENTS
%
% Fichier PARTAGÉ par les quatre versions du document, au même titre que les
% chapitres et que la page de garde : main.tex, main_v2.tex, main_v3.tex et
% main_ensae.tex l'appellent tous. NE PAS LE DUPLIQUER pour faire évoluer une
% version, ce serait quatre textes à maintenir.
%
% UN CHAMP RESTE À COMPLÉTER, ET IL N'APPARTIENT PAS À L'ASSISTANT : le nom de
% famille du fondateur de Nexialog Consulting. Il s'imprime en gras entre
% crochets tant qu'il n'est pas renseigné, donc un dépôt avec un nom absent est
% visible à la relecture plutôt que silencieux.
% =====================================================================

% --- Le champ à compléter --------------------------------------------
\newcommand{\fondateurnexialog}{Ali~BEHBAHANI}
% ---------------------------------------------------------------------

% =====================================================================
%  EXERGUE
%
%  La citation est la « loi de la crédibilité décroissante » de Manski, qui
%  ouvre Identification Problems in the Social Sciences. Elle est à sa place
%  ici et pas seulement par ornement : le mémoire refuse de poser la direction
%  de la contagion et publie un capital en bornes, ce qui est exactement
%  l'arbitrage qu'elle énonce.
%
%  ELLE EST CITÉE DANS SA LANGUE D'ORIGINE DEPUIS LE 28 SEPTEMBRE 2026, demande
%  de Kélian. Elle portait jusque-là une traduction de l'auteur, signalée comme
%  telle ; l'anglais évite d'avoir à défendre une traduction qui n'apporte rien,
%  la phrase étant courte et le lecteur d'un mémoire d'actuariat la lisant sans
%  peine. Pour mémoire, la traduction retirée : « La crédibilité d'une inférence
%  décroît avec la force des hypothèses que l'on maintient. »
%
%  LA CÉSURE EST INTERDITE DANS CE BLOC, et ce n'est pas cosmétique : babel est
%  chargé en français SEUL, donc il couperait cette phrase anglaise sur des
%  motifs français. Interdire la césure ici coûte deux lignes bien réparties et
%  évite de toucher au préambule à deux jours du dépôt.
% =====================================================================
\thispagestyle{empty}
\vspace*{0.30\textheight}

\begin{flushright}
\begin{minipage}{0.72\textwidth}
\itshape\hyphenpenalty=10000 \exhyphenpenalty=10000
The credibility of inference decreases with the strength of the assumptions
maintained.

\vspace{10pt}
\normalfont\small\raggedleft
Charles F. Manski, \emph{Partial Identification of Probability Distributions}, 2003.
\end{minipage}
\end{flushright}

\clearpage

% =====================================================================
%  REMERCIEMENTS
% =====================================================================
\chapter*{Remerciements}
\addcontentsline{toc}{chapter}{Remerciements}
% Voir la note de 20_notes_synthese : un \chapter* laisse le titre courant sur
% la marque precedente. Sans effet visible ici, la page etant unique, mais la
% regle vaut pour les quatre versions et le texte peut s'allonger.
\markboth{Remerciements}{}

% REECRITS LE 1er OCTOBRE 2026 AU SOIR, demande de Kelian : memes personnes, meme ordre, memes fonctions,
% phrases simples, sans formule litteraire ni phrase de cloture. Retires : le paragraphe
% d'ouverture et trois phrases (« Un travail de cette nature... », « Ces quelques mois... »,
% « Son soutien n'a jamais... Ce travail lui revient... »).
Je remercie \fondateurnexialog, fondateur de Nexialog Consulting, de m'avoir accueilli dans
son cabinet et de la confiance qu'il m'a accordée.

Je remercie mon maître de stage, Hugo RAPIOR, responsable R\&D en actuariat, pour son accueil
dans son équipe, sa disponibilité et son encadrement. Nos points d'avancement hebdomadaires et
ses conseils ont beaucoup apporté à ce travail.

Je remercie également les équipes du pôle Actuariat de Nexialog Consulting pour leur accueil et
nos échanges. Je remercie en particulier les trois praticiens qui ont confronté les affirmations
du modèle à leur expérience.

% LES TROIS RELECTEURS NE SONT PAS NOMMÉS, et cette règle ne bouge pas : décision de
% Kélian du 9 septembre 2026, tenue alors même que la première répondante avait coché
% « citez-moi par mon nom ». Ils ne sont désignés ici que par leur nombre, et par leur
% rôle au § de la relecture de praticien de l'annexe C, où leurs réponses sont rapportées.
% Le 28 septembre, la phrase ne citait que la première : elle les cite tous les trois.

Mes remerciements vont aussi au corps enseignant de l'ENSAE Paris et à la filière Actuariat,
pour la qualité de leurs enseignements. Je remercie ma tutrice académique, Caroline HILLAIRET,
pour son suivi et ses remarques sur le traitement de l'incertitude, qui ont fait évoluer ce travail.

Je remercie enfin ma famille pour son soutien.

\clearpage
